\documentclass[10pt,aspectratio=169,mathserif]{beamer}
\usepackage{jsu}
\usepackage[UTF8]{ctex}
\usepackage{amsmath,amsfonts,amssymb,bm}
\usepackage{graphicx,hyperref,url}
\usepackage{booktabs}

\beamertemplateballitem

% ============================================================
% 基本信息（请按需修改）
% ============================================================
\title{演示文稿标题}
\subtitle{副标题（可选）}

\author[姓名]{
    姓名\ 专业班级\\
    学号：xxxxxxxxxx\\
    {\small \url{https://github.com/your-username}}
}

\date{2026年X月}

\begin{document}

%% ------------------------------------------------------------
%% 封面
%% ------------------------------------------------------------
\begin{frame}
    \titlepage
\end{frame}

%% ------------------------------------------------------------
%% 目录
%% ------------------------------------------------------------
\begin{frame}
    \frametitle{目录}
    \tableofcontents
\end{frame}

%% ------------------------------------------------------------
%% 第一部分
%% ------------------------------------------------------------
\section{引言}

\begin{frame}{研究背景}
    \begin{itemize}
        \item 在此处简述研究背景与动机
        \item 指出当前工作中的关键问题
        \item 说明本报告的主要内容与结构
    \end{itemize}

    \vspace{0.5cm}
    \begin{block}{核心问题}
        用一两句话概括你要解决的核心问题。
    \end{block}
\end{frame}

\begin{frame}{报告结构}
    \begin{enumerate}
        \item 引言与问题定义
        \item 方法与技术路线
        \item 实验结果与分析
        \item 总结与展望
    \end{enumerate}
\end{frame}

%% ------------------------------------------------------------
%% 第二部分
%% ------------------------------------------------------------
\section{方法}

\begin{frame}{技术路线}
    \begin{columns}
    \column{0.5\textwidth}
    \textbf{主要步骤：}
    \begin{itemize}
        \item 数据准备与预处理
        \item 模型设计与训练
        \item 评估与对比分析
        \item 结果可视化
    \end{itemize}

    \column{0.5\textwidth}
    \begin{block}{输入}
        原始数据 / 问题描述
    \end{block}
    \begin{exampleblock}{输出}
        结论 / 系统 / 可视化结果
    \end{exampleblock}
    \end{columns}
\end{frame}

\begin{frame}{两栏布局示例}
    \begin{columns}[T]
    \column{0.48\textwidth}
    \textbf{左侧要点}
    \begin{itemize}
        \item 条目 A
        \item 条目 B
        \item 条目 C
    \end{itemize}

    \column{0.48\textwidth}
    \textbf{右侧要点}
    \begin{itemize}
        \item 条目 D
        \item 条目 E
        \item 条目 F
    \end{itemize}
    \end{columns}
\end{frame}

\begin{frame}{公式示例}
    最小二乘目标函数：
    \[
        \min_{\mathbf{w}} \;
        \frac{1}{2n}\sum_{i=1}^{n}
        \bigl(y_i - \mathbf{w}^{\mathsf{T}}\mathbf{x}_i\bigr)^{2}
        + \frac{\lambda}{2}\|\mathbf{w}\|_{2}^{2}
    \]

    \vspace{0.4cm}
    贝叶斯公式：
    \[
        P(\theta\mid D)
        = \frac{P(D\mid\theta)\,P(\theta)}{P(D)}
    \]
\end{frame}

%% ------------------------------------------------------------
%% 第三部分
%% ------------------------------------------------------------
\section{结果}

\begin{frame}{表格示例}
\begin{table}
\centering
\begin{tabular}{lccc}
\toprule
\textbf{方法} & \textbf{指标 A} & \textbf{指标 B} & \textbf{耗时 (s)} \\
\midrule
Baseline     & 0.72 & 0.68 & 12.3 \\
Method A     & 0.81 & 0.79 & 15.6 \\
Method B     & 0.85 & 0.83 & 18.1 \\
\textbf{Ours}& \textbf{0.89} & \textbf{0.87} & 14.2 \\
\bottomrule
\end{tabular}
\end{table}

\vspace{0.4cm}
\begin{center}
\small\textit{表注：示例数据，请替换为实际实验结果。}
\end{center}
\end{frame}

\begin{frame}{强调块示例}
    \begin{block}{普通块（Block）}
        用于说明一般性内容或定义。
    \end{block}

    \begin{alertblock}{警告块（Alert）}
        用于突出需要注意的重要信息。
    \end{alertblock}

    \begin{exampleblock}{示例块（Example）}
        用于展示例子、结果或可复用片段。
    \end{exampleblock}
\end{frame}

\begin{frame}{图片占位}
    \begin{center}
        % 将图片放入 images/ 目录后取消注释：
        % \includegraphics[width=0.7\textwidth]{images/your-figure.png}
        \fbox{\parbox{0.7\textwidth}{\centering\vspace{2.5cm}
        在此插入图片\\[0.3cm]
        {\small\texttt{\textbackslash includegraphics[...]\{images/xxx.png\}}}
        \vspace{2.5cm}}}
    \end{center}
\end{frame}

%% ------------------------------------------------------------
%% 第四部分
%% ------------------------------------------------------------
\section{总结}

\begin{frame}{工作总结}
    \begin{itemize}
        \item 完成了……
        \item 验证了……
        \item 相比基线提升了……
    \end{itemize}

    \vspace{0.5cm}
    \begin{block}{主要贡献}
        \begin{enumerate}
            \item 贡献一
            \item 贡献二
            \item 贡献三
        \end{enumerate}
    \end{block}
\end{frame}

\begin{frame}{展望}
    \begin{itemize}
        \item 扩展到更大规模数据集
        \item 探索更高效的实现方式
        \item 与实际应用场景进一步结合
    \end{itemize}
\end{frame}

%% ------------------------------------------------------------
%% 致谢
%% ------------------------------------------------------------
\section*{}

\begin{frame}
    \vfill
    \centering
    \Huge\textbf{感谢聆听}

    \vspace{1cm}

    \Large 欢迎提问与交流
    \vfill
\end{frame}

\end{document}
